% =============================================================================
% 3. PROBLEM FORMULATION AND ENGINEERING PRELIMINARIES
% =============================================================================
\section{Problem Formulation and Engineering Preliminaries}
\label{sec:problem_formulation}

\subsection{Sequential Edge Few-Shot Task Formulation}
\label{sec:task_formulation}
Field diagnosis requires identifying crop diseases directly on mobile devices as leaves are inspected in the field. When deployed to a new crop or region, the system must calibrate using only a small reference set of labeled leaves collected on site:
\begin{equation}
    \mathcal{S} = \{(\mathbf{x}_i, y_i)\}_{i=1}^{N_s}, \qquad N_s = C \times K,
\end{equation}
where $C$ is the number of target pathology classes (e.g., $C \in \{5, 11, 27\}$) and $K$ is the number of labeled examples per class (typically $K = 5$). The target classes themselves are absent from pre-training ($\mathcal{Y}_{\text{base}} \cap \mathcal{Y}_{\text{novel}} = \emptyset$).

During active scouting, leaves are observed sequentially as an online stream:
\begin{equation}
    \mathcal{Q}_{\text{stream}} = \{\mathbf{x}_t\}_{t=1}^{T}, \qquad \mathbf{x}_t \in \mathbb{R}^{H \times W \times 3}.
\end{equation}
The system must output an immediate diagnosis $\hat{y}_t \in \{1, \dots, C\}$ for each incoming frame $\mathbf{x}_t$, without buffering queries or waiting for future images.

Performing this diagnosis on edge hardware faces strict computational limits. Fine-tuning deep backbones on device requires substantial memory overhead for activation caching and optimizer state, causing severe processor latency~\cite{lin2022device,radocaj2025evolution}. Our model therefore relies on frozen foundation encoders ($\nabla_\theta = \emptyset$). We extract complementary visual and biological features using DINOv2-Small ($d_1 = 384$) and BioCLIP ($d_2 = 512$), yielding a fixed joint descriptor $\mathbf{z} \in \mathbb{R}^d$ ($d = 896$) without gradient calculations during field adaptation.

\subsection{Engineering and Operational Constraints}
\label{sec:engineering_constraints}
Because agricultural sites often have limited or unreliable connectivity, foliar diagnosis should run on edge devices rather than depend on cloud servers~\cite{radocaj2025evolution,sharif2026real}. We define the on-device diagnostic objective as a constrained optimization problem that maximizes accuracy while bounding the disease escape rate (false-negative rate, FNR) and satisfying operational engineering boundaries:
\begin{equation}
\begin{aligned}
    \max_{\mathbf{W}, \mathbf{b}, \, \tau_{\text{conf}}} \quad & \mathbb{E}_{(\mathbf{x}_t, y_t)}\big[\mathbb{I}(\hat{y}_t = y_t)\big] \\
    \text{subject to} \quad & \begin{cases}
        \text{FNR}(\hat{y}) \le \epsilon_{\text{esc}} & \text{(agronomic safety)} \\
        \overline{L}_{\text{inf}}(\tau_{\text{conf}}) \le L_{\max} & \text{(real-time stream latency)} \\
        T_{\text{adapt}}(\mathcal{S}) \le T_{\max}, \; \nabla_\theta = \emptyset & \text{(gradient-free calibration)} \\
        \text{Mem}_{\text{PSS}} \le M_{\max} & \text{(mobile thermal and memory budget)},
    \end{cases}
\end{aligned}
\label{eq:constrained_objective}
\end{equation}
where $\mathbf{W} \in \mathbb{R}^{C \times d}$ and $\mathbf{b} \in \mathbb{R}^C$ denote the analytical projection parameters, and $\tau_{\text{conf}}$ governs the two-tier escalation gate.

While the exact operational parameters $(\epsilon_{\text{esc}}, L_{\max}, T_{\max}, M_{\max})$ depend on the scouting platform and target crop pathology, we establish representative target boundaries for our edge deployment:
\begin{enumerate}
    \item \textbf{Agronomic Disease Escape Limit ($\text{FNR} \le \epsilon_{\text{esc}}$):} In plant protection, misclassifying an infected leaf as healthy lets pathogens spread unnoticed. Integrated pest management treats missed infections as the costliest error~\cite{richard2022control}, so we set our own target at an FNR below $\epsilon_{\text{esc}} = 1.0\%$.
    \item \textbf{Real-Time Stream Latency ($\overline{L}_{\text{inf}} \le L_{\max}$):} Field scouting platforms scan crop canopies continuously. At the 1--2\,m/s row speeds typical of field scouting, 6\,fps corresponds to a 0.17--0.33\,m stride, roughly one frame per leaf cluster. This sets a 170\,ms per-frame budget ($1/6$\,s $\approx 167$\,ms, rounded up), which we adopt as our design requirement. Staying within it on device avoids the latency and bandwidth cost of off-board processing~\cite{radocaj2025evolution}.
    \item \textbf{Sub-Second On-Device Calibration ($T_{\text{adapt}} \le T_{\max}$):} When encountering a new regional pathogen variant in the field, calibrating the model with a few reference leaves ($K=5$) must finish within sub-second time ($T_{\max} \le 1.0$\,s) so scouting can resume immediately. In addition, backpropagation and optimizer tracking are disallowed on device ($\nabla_\theta = \emptyset$) to prevent memory bloat and processor slowdown.
    \item \textbf{Memory Footprint ($\text{Mem}_{\text{PSS}} \le M_{\max}$):} Mobile edge terminals share limited random-access memory (RAM) across the operating system and camera drivers. Setting a budget of $M_{\max} = 500$\,MB peak PSS ensures the diagnostic engine avoids memory-pressure termination and operates safely on fanless embedded hardware.
\end{enumerate}

\subsection{Small-Sample Covariance Degeneracy}
\label{sec:covariance_degeneracy}
Foliar representations exhibit non-uniform feature variance across latent dimensions. Unrelated background patterns---such as green pigmentation, leaf veins, and surface glare---span the highest-variance directions in the feature space. When classification relies on isotropic metrics such as unweighted Euclidean distance or cosine similarity:
\begin{equation}
    \hat{y}_t = \arg\max_{c \in \{1, \dots, C\}} \frac{\mathbf{z}_t^T \boldsymbol{\mu}_c}{\|\mathbf{z}_t\|_2 \|\boldsymbol{\mu}_c\|_2}, \qquad \boldsymbol{\mu}_c = \frac{1}{|\mathcal{S}_c|}\sum_{\mathbf{x} \in \mathcal{S}_c} \mathbf{z}(\mathbf{x}),
    \label{eq:isotropic_cosine}
\end{equation}
all feature coordinates are weighted equally ($\mathbf{\Sigma} = \sigma^2 \mathbf{I}$). As a result, distance metrics are dominated by healthy leaf background rather than subtle lesion symptoms.

Fisher Linear Discriminant Analysis resolves this issue by projecting features through the inverse within-class covariance $\hat{\mathbf{\Sigma}}^{-1}$, which suppresses high-variance shared background directions and amplifies discriminative lesion signals. Computing this projection requires the empirical pooled within-class scatter matrix:
\begin{equation}
    \mathbf{S}_W = \frac{1}{N - C} \sum_{c=1}^C \sum_{\mathbf{z} \in \mathcal{S}_c} (\mathbf{z} - \boldsymbol{\mu}_c)(\mathbf{z} - \boldsymbol{\mu}_c)^T \in \mathbb{R}^{d \times d}.
    \label{eq:empirical_scatter}
\end{equation}
In field deployment, however, the number of support samples $N$ is far smaller than the feature dimension $d$ ($N \ll d$). Across our target evaluation settings:
\begin{itemize}
    \item In 5-way 5-shot field diagnosis, $N = 25$ versus $d = 896$, yielding $\text{rank}(\mathbf{S}_W) \le 20$.
    \item On the 11-class PlantVillage benchmark without augmentation, $N = 55$ versus $d = 896$, yielding $\text{rank}(\mathbf{S}_W) \le 44$.
    \item Even after multi-view support augmentation ($N = 440$), $\text{rank}(\mathbf{S}_W) \le 429 < 896$.
\end{itemize}
Because $\text{rank}(\mathbf{S}_W) \le N - C \ll d$, the empirical scatter matrix $\mathbf{S}_W$ has nullity $d - (N - C) > 0$. This makes $\mathbf{S}_W$ strictly singular ($\det(\mathbf{S}_W) = 0$) and numerically unstable under simple diagonal loading. This fundamental singularity motivates the analytical shrinkage and uncertainty-gated architecture detailed in Section~\ref{sec:methodology}.
